\documentclass[11pt]{article}
\usepackage{amsmath,amssymb,amsthm}
\usepackage[margin=1in]{geometry}
\newtheorem{theorem}{Theorem}
\newtheorem{lemma}{Lemma}
\newcommand{\F}{\mathbb F_2}
\title{Spine tree--DAG separation for affine-cell branching}
\author{NC0\_k-Avoid Laboratory, V68}
\date{2026-07-31}
\begin{document}
\maketitle

\paragraph{Convention.}
For local coordinates $(a,b,c)$, truth-table index means $4a+2b+c$.
The positive fiber of mask $\mathtt{0x07}$ is
$\{000,001,010\}$ and therefore pins $a=0$.

\paragraph{Construction.}
For $k\geq 1$, use variables
$s,(u_0,v_0),\ldots,(u_{k-1},v_{k-1})$.
For every $t$, include two motif gates on $(s,u_t,v_t)$ and
$(s,v_t,u_t)$ with affine cells
\[
 C^0_t=\{000\},\qquad C^1_t=\{001,010\}.
\]
On motif zero add two anchors with positive fibers
$\{000,001,011\}$ and $\{000,010,011\}$, each partitioned into
$\{000\}$ and the complementary two-point affine line.
The anchor masks are $\mathtt{0x0b}$ and $\mathtt{0x0d}$, both in the NPN
orbit of $\mathtt{0x07}$.
Thus $n=2k+1$ and $m=2k+2=n+1$.

\begin{theorem}[Spine tree--DAG separation]
Let $S_k$ be the construction above.
The number of consistent complete branch signatures is
\[
 c(S_k)=2^{k-1}=2^{(n-3)/2}.
\]
Hence every complete inconsistency-pruned affine-cell branching tree for
$S_k$ has at least $2^{k-1}$ leaves.
Under the fixed gate order described in the proof and existential projection
onto variables occurring in remaining gates, there is a residual-state DAG
with exactly $3k+4$ nonterminal states, plus accepting and inconsistent
terminals.
\end{theorem}

\begin{proof}
Every cell forces $s=0$.  For each $t\geq1$, the assignment
$(u_t,v_t)=(0,0)$ gives motif branch pair $00$, while the two one-hot
assignments give branch pair $11$; $(1,1)$ is outside the motif fiber.
Thus every free motif contributes exactly two signatures.

For $t=0$, the $\mathtt{0x0b}$ anchor excludes $(1,0)$, the
$\mathtt{0x0d}$ anchor excludes $(0,1)$, and the motif excludes $(1,1)$.
Only $(0,0)$ remains.  Conditioned on $s=0$, the other pairs are disjoint,
so their choices multiply and give $2^{k-1}$ signatures.  Distinct complete
signatures require distinct complete consistent leaves.

Order the gates as the two motif-zero gates, the two anchors, and then the two
gates of each free motif consecutively.  The prefix has seven nonterminal
projected residual states.  In a free motif, the first gate creates the two
affine residuals $u_t=v_t=0$ and $u_t+v_t=1$.  The duplicate second gate is
forced in each residual.  After it is consumed, $u_t,v_t$ occur nowhere in
the suffix; existential projection removes their equations and both paths
merge back to the same constraint $s=0$.  Each free motif adds three states,
so the count is $7+3(k-1)=3k+4$.
\end{proof}

\paragraph{Boundary.}
The projected quantity used here is repository-local.  This theorem does not identify it with OBDD size, FBDD size, resolution, Res-Lin, or any other standard proof-complexity measure.  It does not prove a general polynomial DAG bound, unrestricted $\mathrm{NC}^0_3$-Avoid, a circuit lower bound, or a separation of P from NP.

\end{document}
